\begin{proof}[Proof of \cref{obj:thm:observed-record-precision-frontier}]
\begin{enumerate}
\item \textit{Upper risk bound.}
Fix admissible \(n,d,q\). Define the association-revelation probability and a numerical constant by
\[
p=1-(1-q)^d,
\qquad
\alpha_{\mathrm{ret}}=1-e^{-1}\in(0,1).
\]
We first establish
\[
\alpha_{\mathrm{ret}}\min\{1,dq\}\le p\le\min\{1,dq\}.
\]
For the upper bound, induction on the nonnegative integer degree starts with
\(1-(1-q)^0=0\). The induction step is
\[
\begin{aligned}
1-(1-q)^{d+1}
&=q+(1-q)\bigl[1-(1-q)^d\bigr]\\
&\le q+(1-q)dq
=(d+1)q-dq^2
\le(d+1)q.
\end{aligned}
\]
Also \(p\le1\), since \(1-q\ge0\).

For the lower bound, the exponential inequality \(1-q\le e^{-q}\) gives
\[
(1-q)^d\le e^{-dq},
\qquad
p\ge1-e^{-dq}.
\]
If \(dq\le1\), convexity of the exponential yields
\[
e^{-dq}\le(1-dq)e^0+dq\,e^{-1},
\qquad
1-e^{-dq}\ge\alpha_{\mathrm{ret}}dq.
\]
If \(dq>1\), monotonicity gives
\[
e^{-dq}\le e^{-1},
\qquad
1-e^{-dq}\ge\alpha_{\mathrm{ret}}.
\]
In particular, \(p>0\) whenever \(q>0\).

By \cref{obj:thm:observable-upper,obj:def:frontier-rule},
\[
R_{n,d}(q)
\le
\sup_{\theta\in\mathcal M_{n,d}}
\mathcal L(T^\star_{n,d,q};\theta,q)
\le\min\{1,u(n,d,q)\},
\]
and the attaining rule is measurable and real-valued. For \(q>0\), all denominators below are positive, and the probability bounds imply
\[
\frac{d^2}{n}\le\frac{d^2}{np},
\qquad
\frac{d}{nq}\le\frac{d^2}{np}.
\]
Using \(d+1\le2d\) and \(1-q\le1\) in
\cref{obj:def:upper-envelope}, we obtain
\[
\begin{aligned}
u(n,d,q)
&=\frac{5(d+1)^2}{n}+\frac{4d(1-q)}{nq}\\
&\le20\frac{d^2}{n}+4\frac{d}{nq}
\le24\frac{d^2}{np}.
\end{aligned}
\]
If \(d^2/(np)\le1\), then
\[
\min\{1,u(n,d,q)\}
\le u(n,d,q)
\le24\frac{d^2}{np}
=24F(n,d,q).
\]
If \(d^2/(np)>1\), then
\[
\min\{1,u(n,d,q)\}\le1\le24=24F(n,d,q).
\]
At \(q=0\), the envelope is \(+\infty\) and the scale is \(1\), so
\[
\min\{1,u(n,d,0)\}=1\le24F(n,d,0).
\]
Consequently, throughout the admissible range,
\[
R_{n,d}(q)
\le
\sup_{\theta\in\mathcal M_{n,d}}
\mathcal L(T^\star_{n,d,q};\theta,q)
\le24F(n,d,q).
\]
% lean: CausalSmith.Experimentation.ThinnedgraphAdditiveRiskFrontier.frontierEstimator_risk_upper

\item \textit{Lower risk bound.}
The definition gives \(F(n,d,q)\le1\). First suppose that
\(d^2/n\ge1/16\). The canonical product design satisfies
\cref{obj:ass:assignment,obj:ass:audit,obj:ass:design-independence}, so
\cref{obj:lem:supplied-block-record-lower} supplies
\[
\begin{aligned}
R_{n,d}(q)
&\ge\frac{1}{160000\pi^2}
       \min\left\{1,\frac{(d+1)^2}{n}\right\}\\
&\ge\frac{1}{2560000\pi^2}
\ge\frac{F(n,d,q)}{2560000\pi^2}.
\end{aligned}
\]
Here the second inequality follows from
\[
\frac{(d+1)^2}{n}\ge\frac{d^2}{n}\ge\frac1{16}.
\]

Now suppose that \(d^2/n<1/16\). Since \(d\ge1\),
\[
4d\le16d^2<n.
\]
Define the block count and source count by
\[
B=\left\lfloor\frac{n}{2d}\right\rfloor,
\qquad
m=Bd.
\]
The floor inequalities give
\[
2m\le n<2m+2d.
\]
Together with \(4d<n\), these imply
\[
B\ge1,
\qquad
\frac n4\le m\le\frac n2,
\qquad
2Bd\le n.
\]
Use the fixed source and recipient construction of \cref{obj:synth_4}.

Define the testing scale, numerical constant, and chosen amplitude by
\[
s=
\begin{cases}
\min\{1,d/\sqrt{mp}\},&p>0,\\
1,&p=0,
\end{cases}
\qquad
\kappa=\frac1{100\pi},
\qquad
h=\kappa s.
\]
Since \(m,d>0\), the definition gives \(0<s\le1\). Moreover,
\[
0<h\le\kappa<\frac14,
\]
where the last inequality follows from \(\pi>3\).

We next compare the risk and testing scales:
\[
F(n,d,q)\le s^2.
\]
For \(q>0\), we have \(p>0\) and
\[
\left(\frac d{\sqrt{mp}}\right)^2=\frac{d^2}{mp},
\qquad
\frac{d^2}{np}\le\frac{d^2}{mp}.
\]
If \(d/\sqrt{mp}\le1\), then
\[
F(n,d,q)\le\frac{d^2}{np}\le\frac{d^2}{mp}=s^2.
\]
If \(d/\sqrt{mp}>1\), then \(s^2=1\ge F(n,d,q)\).
For \(q=0\), \(p=0\) and \(F(n,d,0)=s^2=1\).
Thus the comparison covers the entire retention range.

Take the opposite-sign priors of \cref{obj:def:block-prior} at amplitude \(h\).
By \cref{obj:lem:block-law}, they are supported almost surely in
\(\mathcal M_{n,d}\), with constant targets \(\pm\Delta\), where
\[
\Delta=\frac{mh}{n}>0.
\]
The record maps are jointly measurable: the partition and binary coordinates
are finite, and the outcomes depend affinely on the real baseline coordinates.
The block and amplitude conditions established above allow
\cref{obj:thm:block-testing-scale} to give
\[
\operatorname{TV}(P_{+1,h},P_{-1,h})\le\frac12.
\]
Applying \cref{obj:lem:full-record-two-prior-risk} therefore yields
\[
R_{n,d}(q)
\ge\frac{\Delta^2}{2}
       \bigl(1-\operatorname{TV}(P_{+1,h},P_{-1,h})\bigr)
\ge\frac{\Delta^2}{4}.
\]
Finally, \(m/n\ge1/4\) implies
\[
\frac{\kappa s}{4}
\le\frac{m\kappa s}{n}
=\Delta.
\]
Hence
\[
\begin{aligned}
\frac{F(n,d,q)}{2560000\pi^2}
&\le\frac{s^2}{2560000\pi^2}
\le\frac{s^2}{640000\pi^2}\\
&=\frac14\left(\frac{\kappa s}{4}\right)^2
\le\frac{\Delta^2}{4}
\le R_{n,d}(q).
\end{aligned}
\]
This proves the lower bound in both degree regimes, including \(q=0\).
% lean: CausalSmith.Experimentation.ThinnedgraphAdditiveRiskFrontier.frontier_minimax_lower

\item \textit{Necessity of the sequence conditions.}
Fix admissible sequences \(d_n,q_n\), and suppose that measurable randomized
estimators \(T_n\) satisfy
\[
\rho_{\mathrm{risk},n}
:=
\sup_{\theta\in\mathcal M_{n,d_n}}
\mathcal L(T_n;\theta,q_n)
\longrightarrow0,
\qquad n\ge4.
\]
The preceding lower bound and the definition of minimax risk give
\[
0\le
\frac{F(n,d_n,q_n)}{2560000\pi^2}
\le R_{n,d_n}(q_n)
\le\rho_{\mathrm{risk},n}.
\]
Squeezing to zero and multiplying by the positive constant
\(2560000\pi^2\), we obtain
\[
F(n,d_n,q_n)\longrightarrow0.
\]

To extract the sequence conditions, we establish the endpoint and
positive-retention comparisons used here. Directly from
\cref{obj:def:frontier-scale}, since \(d\ge1\),
\[
F(n,d,0)=1,
\qquad
F(n,d,1)=\min\left\{1,\frac{d^2}{n}\right\}.
\]
For \(0<q\le1\), the established bounds on \(p\) give
\[
\frac{d^2}{n}\le\frac{d^2}{np},
\qquad
\frac{d}{nq}\le\frac{d^2}{np}.
\]
The lower probability estimate also gives, by splitting at \(dq=1\),
\[
\alpha_{\mathrm{ret}}\frac{d^2}{np}
\le
\begin{cases}
\dfrac{d}{nq},&dq\le1,\\[4pt]
\dfrac{d^2}{n},&dq>1.
\end{cases}
\]
Consequently,
\[
\frac{d^2}{n}+\frac{d}{nq}\le2\frac{d^2}{np},
\qquad
\alpha_{\mathrm{ret}}\frac{d^2}{np}
\le\frac{d^2}{n}+\frac{d}{nq}.
\]

For the lower truncated comparison, if \(d^2/(np)\ge1\), then
\[
\frac12\min\left\{1,\frac{d^2}{n}+\frac{d}{nq}\right\}
\le1=F(n,d,q).
\]
If \(d^2/(np)<1\), then
\[
\frac12\min\left\{1,\frac{d^2}{n}+\frac{d}{nq}\right\}
\le\frac12\left(\frac{d^2}{n}+\frac{d}{nq}\right)
\le\frac{d^2}{np}
=F(n,d,q).
\]
For the upper comparison, \(0<\alpha_{\mathrm{ret}}<1\) gives both
\[
\alpha_{\mathrm{ret}}F(n,d,q)\le1,
\qquad
\alpha_{\mathrm{ret}}F(n,d,q)
\le\alpha_{\mathrm{ret}}\frac{d^2}{np}
\le\frac{d^2}{n}+\frac{d}{nq}.
\]
Combining these inequalities proves
\[
\frac12\min\left\{1,\frac{d^2}{n}+\frac{d}{nq}\right\}
\le F(n,d,q)
\le
\alpha_{\mathrm{ret}}^{-1}
\min\left\{1,\frac{d^2}{n}+\frac{d}{nq}\right\}.
\]
% lean: CausalSmith.Experimentation.ThinnedgraphAdditiveRiskFrontier.frontier_elbow

Since \(F(n,d_n,q_n)\to0\), eventually
\[
F(n,d_n,q_n)<\frac12.
\]
The zero-retention identity then forces \(q_n>0\) eventually. On that tail,
the lower comparison implies
\[
\frac{d_n^2}{n}+\frac{d_n}{nq_n}<1:
\]
otherwise its truncated value would be \(1\), forcing
\(F(n,d_n,q_n)\ge1/2\). Removing the truncation now yields
\[
0\le\frac{d_n^2}{n}
\le\frac{d_n^2}{n}+\frac{d_n}{nq_n}
\le2F(n,d_n,q_n),
\]
and
\[
0\le\frac{d_n}{nq_n}
\le\frac{d_n^2}{n}+\frac{d_n}{nq_n}
\le2F(n,d_n,q_n).
\]
Both real ratios therefore tend to zero.

Define the extended audit ratio for every \(n\ge4\) by
\[
\rho_{\mathrm{audit},n}
=
\begin{cases}
+\infty,&q_n=0,\\[2pt]
\dfrac{d_n}{nq_n},&q_n>0.
\end{cases}
\]
It agrees eventually with the second real ratio, so
\[
\frac{d_n^2}{n}\longrightarrow0,
\qquad
\rho_{\mathrm{audit},n}\longrightarrow0.
\]
% lean: CausalSmith.Experimentation.ThinnedgraphAdditiveRiskFrontier.frontierScale_tendsto_iff

\item \textit{Sufficiency of the sequence conditions.}
Conversely, suppose
\[
\frac{d_n^2}{n}\longrightarrow0,
\qquad
\rho_{\mathrm{audit},n}\longrightarrow0.
\]
Convergence of the extended ratio gives
\(\rho_{\mathrm{audit},n}<1\) eventually. Thus \(q_n>0\) eventually,
and on that tail its finite real value satisfies
\[
\frac{d_n}{nq_n}\longrightarrow0.
\]
It follows that
\[
\frac{d_n^2}{n}+\frac{d_n}{nq_n}\longrightarrow0.
\]
The upper scale comparison proved above gives, eventually,
\[
0\le F(n,d_n,q_n)
\le\alpha_{\mathrm{ret}}^{-1}
\left(\frac{d_n^2}{n}+\frac{d_n}{nq_n}\right),
\]
and hence \(F(n,d_n,q_n)\to0\).

Choose the measurable rule sequence
\[
T_n(O,\xi)=T^\star_{n,d_n,q_n}(O),
\qquad n\ge4.
\]
The upper risk bound then supplies
\[
0\le
\sup_{\theta\in\mathcal M_{n,d_n}}
\mathcal L(T_n;\theta,q_n)
\le24F(n,d_n,q_n)
\longrightarrow0.
\]
Together with necessity, this proves the asserted equivalence.
% lean: CausalSmith.Experimentation.ThinnedgraphAdditiveRiskFrontier.precision_frontier_explicit

\item \textit{Universal constants and endpoint scales.}
Set
\[
c=\frac1{2560000\pi^2},
\qquad
C_0=24,
\qquad
c_1=\frac12,
\qquad
C_1=(1-e^{-1})^{-1}.
\]
These constants are positive because \(\pi>0\) and \(e^{-1}<1\).
The upper and lower risk bounds establish the risk chain with \(c,C_0\);
the sequence argument establishes consistency in both directions.
The endpoint identities and positive-retention comparison established above
give the remaining assertions with \(c_1,C_1\).
% lean: CausalSmith.Experimentation.ThinnedgraphAdditiveRiskFrontier.observed_record_precision_frontier
\end{enumerate}
\end{proof}
